% !TeX root = ../main.tex
\section{Data Training Set}
\label{sec:training_set}
For the data training set, a simulator was used. In a previous phase of the work, sensors were simulated to send data based on physical models. This section describes the physical variables used to generate simulated telemetry values and the relations among them. \textbf{Real-world data may also be incorporated to build a mixed training dataset composed of real and simulated measurements}.

\subsection{Telemetry Metric Fault Isolation Framework}

Autonomous Mobile Robots (AMRs) and Automated Guided Vehicles (AGVs) rely on high-frequency telemetry streams to monitor powertrain health. Telemetry anomalies are categorized into univariate sensor-level faults and multivariate system-level electromechanical failures \cite{Pillai2024Taxonomy}.

The multivariate fault rules described in this section are supported by a simplified electromechanical model of the actuation and energy subsystems. This model is used in the simulator to generate physically plausible correlations among telemetry channels, and later in the diagnostic rules to interpret anomalous combinations of measurements.

\subsection{Univariate Sensor Anomalies}
\begin{itemize}
    \item \textbf{Motor Current ($I_{\mathrm{m}}$):} Isolated high-frequency noise or constant offset shift. Caused by shunt resistor thermal aging or Hall-effect sensor gain drift without mechanical load changes \cite{Pillai2024Taxonomy}.
    \item \textbf{Motor Voltage ($V_{\mathrm{m}}$):} Telemetry dropouts or step-discontinuities. Caused by analog voltage divider resistor degradation or microcontroller ADC channel failure \cite{Kim2022HealthMonitoring}.
    \item \textbf{Motor Temperature ($T_{\mathrm{m}}$):} Out-of-range or saturated readings, represented as $T_{\mathrm{m}} \to T_{\mathrm{max,ADC}}$ or $T_{\mathrm{m}} \to T_{\mathrm{min,ADC}}$. Caused by NTC/PT100 sensor open-circuit, signal line short-circuit, or thermal interface material (TIM) degradation \cite{Li2023Prognostics, Pillai2024Taxonomy}.
    \item \textbf{Linear Velocity ($v$):} Pulse dropouts or instantaneous zero velocity readings during active motion. Caused by optical encoder disk contamination, quadrature channel loss, or magnetic pickup misalignment \cite{Kim2022HealthMonitoring}.
    \item \textbf{Battery State ($\mathrm{SOC}, V_{\mathrm{bat}}$):} Instantaneous SOC percentage jumps or false charging flags. Caused by Coulomb-counting integration drift or CAN-bus telemetry frame corruption \cite{Kim2022HealthMonitoring}.
\end{itemize}

\subsection{Multivariate Fault Diagnostic Matrix}

The interaction of current, voltage, temperature, velocity, and battery metrics under complex failure modes is governed by electromechanical differential equations \cite{Todeschini2015Actuator}.

The electrical behavior of the motor is represented by the armature-equivalent circuit:
\begin{align}
    e_{\mathrm{b}}(t) &= K_{\mathrm{e}} \omega_{\mathrm{m}}(t) \label{eq:motor_back_emf} \\
    V_{\mathrm{m}}(t) &= R_{\mathrm{a}}(T_{\mathrm{m}}(t)) I_{\mathrm{m}}(t) + L_{\mathrm{a}} \dot{I}_{\mathrm{m}}(t) + e_{\mathrm{b}}(t) \label{eq:motor_voltage}
\end{align}

The mechanical behavior of the actuator is represented by the electromagnetic torque and rotor dynamics equations:
\begin{align}
    \tau_{\mathrm{m}}(t) &= K_{\mathrm{t}} I_{\mathrm{m}}(t) \label{eq:electromagnetic_torque} \\
    J_{\mathrm{eq}} \dot{\omega}_{\mathrm{m}}(t) &= \tau_{\mathrm{m}}(t) - B_{\mathrm{m}} \omega_{\mathrm{m}}(t) - \tau_{\mathrm{L,ref}}(t) \label{eq:torque_balance} \\
    \tau_{\mathrm{L,ref}}(t) &= \frac{\tau_{\mathrm{L}}(t)}{\eta_{\mathrm{g}} i_{\mathrm{g}}} \label{eq:reflected_load_torque}
\end{align}

The battery terminal voltage is represented by a simplified equivalent-circuit model:
\begin{equation}
    V_{\mathrm{bat}}(t) = V_{\mathrm{ocv}}(\mathrm{SOC}(t)) - I_{\mathrm{bat}}(t) R_{\mathrm{int}}(\mathrm{SOC}(t), T_{\mathrm{bat}}(t)) - V_{\mathrm{pol}}(t) \label{eq:battery_ecm}
\end{equation}

The main variables and units are:
\begin{description}
    \item[$V_{\mathrm{m}}(t)$] Motor terminal voltage in volts ($\mathrm{V}$). This is the voltage measured or simulated at the motor equivalent terminals.
    \item[$I_{\mathrm{m}}(t)$] Armature current or torque-producing current in amperes ($\mathrm{A}$). It is one of the main indicators of mechanical load and electrical faults.
    \item[$T_{\mathrm{m}}(t)$] Motor or winding temperature in degrees Celsius ($^\circ\mathrm{C}$). It affects copper losses and, indirectly, winding resistance.
    \item[$R_{\mathrm{a}}(T_{\mathrm{m}}(t))$] Armature or equivalent winding resistance in ohms ($\Omega$). A common thermal approximation is $R_{\mathrm{a}}(T_{\mathrm{m}}) \approx R_{\mathrm{a,ref}}[1+\alpha_{\mathrm{Cu}}(T_{\mathrm{m}}-T_{\mathrm{ref}})]$.
    \item[$L_{\mathrm{a}}$] Armature inductance in henries ($\mathrm{H}$). It mainly affects transient current changes.
    \item[$e_{\mathrm{b}}(t)$] Back-electromotive force in volts ($\mathrm{V}$). It opposes the applied voltage and is proportional to motor speed.
    \item[$K_{\mathrm{e}}$] Back-EMF constant in volts per radian per second ($\mathrm{V}/(\mathrm{rad}/\mathrm{s})$).
    \item[$\omega_{\mathrm{m}}(t)$] Motor angular velocity in radians per second ($\mathrm{rad}/\mathrm{s}$). Under no-slip conditions, linear velocity can be related to motor speed by $v(t) = r_{\mathrm{w}} \omega_{\mathrm{m}}(t) / i_{\mathrm{g}}$, where $r_{\mathrm{w}}$ is the wheel radius and $i_{\mathrm{g}}$ is the gearbox ratio.
    \item[$\tau_{\mathrm{m}}(t)$] Electromagnetic torque in newton-meters ($\mathrm{N}\,\mathrm{m}$).
    \item[$K_{\mathrm{t}}$] Torque constant in newton-meters per ampere ($\mathrm{N}\,\mathrm{m}/\mathrm{A}$).
    \item[$J_{\mathrm{eq}}$] Equivalent inertia referred to the motor shaft in kilograms times square meters ($\mathrm{kg}\,\mathrm{m}^2$). If motor and load inertias are included, it can be written as $J_{\mathrm{eq}} = J_{\mathrm{m}} + J_{\mathrm{L}}/i_{\mathrm{g}}^2$, neglecting gearbox inertia.
    \item[$B_{\mathrm{m}}$] Viscous friction coefficient in newton-meters per radian per second ($\mathrm{N}\,\mathrm{m}/(\mathrm{rad}/\mathrm{s})$).
    \item[$\tau_{\mathrm{L}}(t)$] External load torque measured at the load side in newton-meters ($\mathrm{N}\,\mathrm{m}$).
    \item[$\tau_{\mathrm{L,ref}}(t)$] Load torque referred to the motor shaft in newton-meters ($\mathrm{N}\,\mathrm{m}$). The gearbox efficiency $\eta_{\mathrm{g}}$ and ratio $i_{\mathrm{g}}$ are used in this transformation.
    \item[$V_{\mathrm{bat}}(t)$] Battery terminal voltage under load in volts ($\mathrm{V}$).
    \item[$V_{\mathrm{ocv}}(\mathrm{SOC}(t))$] Open-circuit voltage of the battery in volts ($\mathrm{V}$). It depends on the state of charge.
    \item[$\mathrm{SOC}(t)$] State of charge, represented as a dimensionless quantity in $[0,1]$ when used in equations. If a percentage is shown, it should be explicitly written as $\mathrm{SOC}_{\%}$ to avoid ambiguity.
    \item[$I_{\mathrm{bat}}(t)$] Battery current in amperes ($\mathrm{A}$). Its sign convention should be defined by the simulator: discharge is positive in the equation used here.
    \item[$R_{\mathrm{int}}(\mathrm{SOC},T_{\mathrm{bat}})$] Battery internal resistance in ohms ($\Omega$). It depends on state of charge and temperature, and it usually increases with state of health degradation.
    \item[$V_{\mathrm{pol}}(t)$] Polarization voltage in volts ($\mathrm{V}$). It groups transient electrochemical effects, such as diffusion and activation overpotentials, that are not represented by pure ohmic drop.
\end{description}
For computational efficiency and parameter tractability, $V_{\mathrm{pol}}(t)$ is assumed to be zero ($V_{\mathrm{pol}}(t) = 0$), reducing the energy storage model to a pure ohmic equivalent circuit ($V_{\mathrm{bat}}(t) = V_{\mathrm{ocv}}(\mathrm{SOC}) - I_{\mathrm{bat}}(t) \cdot R_{\mathrm{int}}$). This simplification is physically justified because the primary objective of this simulation is high-frequency electromechanical fault detection (e.g., motor stalls, wheel slippage, bearing wear), where instantaneous ohmic voltage sag dominates. Furthermore, static ohmic resistance $R_{\mathrm{int}}$ can be directly parameterized from manufacturer datasheets without requiring complex Electrochemical Impedance Spectroscopy (EIS) testing.


These equations justify why some anomalies cannot be detected by looking at a single signal in isolation: a mechanical fault can modify current, speed, temperature, and battery voltage simultaneously. The following table translates those physical relations into engineering rules for the simulated telemetry. The symbols $\uparrow$ and $\downarrow$ indicate deviation from a healthy operating profile, and labels such as \textit{Nominal}, \textit{Low}, \textit{Base}, and \textit{Baseline} are configurable thresholds computed from simulator profiles or historical fleet data.

\begin{table}[htbp]
\centering
\caption{Multivariate Telemetry Diagnostic Rules for AGV/AMR Powertrains}
\label{tab:diagnostic_matrix}
\small
\begin{tabular}{p{2.8cm}p{3.5cm}p{4.2cm}p{4.0cm}}
\toprule
\textbf{Fault State} & \textbf{Telemetry Pattern} & \textbf{Physical Mechanism} & \textbf{Diagnostic Condition Rule} \\
\midrule
\textbf{Mechanical Stall / Rotor Lock} & $I_{\mathrm{m}} \to I_{\mathrm{stall}}$, $v = 0$, $V_{\mathrm{m}} \to V_{\mathrm{bat}}$, $\frac{\mathrm{d}T_{\mathrm{m}}}{\mathrm{d}t} \gg 0$, $V_{\mathrm{bat}} \downarrow\downarrow$ & Mechanical blockage prevents rotor rotation ($\omega_{\mathrm{m}}=0$), back-EMF collapses ($E_{\mathrm{b}}=0$), causing pure resistive current surge \cite{Pillai2024Taxonomy, Li2023Prognostics}. & $(I_{\mathrm{m}} > 2.5 \cdot I_{\mathrm{rated}}) \land (v < 0.01) \land (v_{\mathrm{ref}} > 0.1)$ \\
\addlinespace
\textbf{Wheel Slippage} & $v_{\mathrm{enc}} \gg v_{\mathrm{ground}}$, $I_{\mathrm{m}} \downarrow$, $V_{\mathrm{m}} = \text{Nominal}$, $T_{\mathrm{m}} = \text{Low}$ & Friction coefficient drops ($\mu \to 0$), load torque collapses ($\tau_{\mathrm{L}} \to 0$), motor accelerates with low current \cite{Reina2009Doppler}. & $\left(\frac{v_{\mathrm{enc}} - v_{\mathrm{ground}}}{v_{\mathrm{ground}}} > 0.15\right) \land (I_{\mathrm{m}} < I_{\mathrm{base}})$ \\
\addlinespace
\textbf{Gearbox / Bearing Wear} & $I_{\mathrm{m}} \uparrow \text{ (baseline)}$, $v = v_{\mathrm{ref}}$, $T_{\mathrm{m,ss}} \uparrow$, $V_{\mathrm{bat}} = \text{Nominal}$ & Bearing race pitting or gear misalignment increases mechanical friction torque ($\tau_{\mathrm{friction}}$) \cite{Li2023Prognostics, Gultekin2022SignalFusion}. & $(\Delta I_{\mathrm{m,\text{mean}}} > +20\%) \land (v = v_{\mathrm{ref}})$ \\
\addlinespace
\textbf{Stator Inter-Turn Short} & $I_{\mathrm{m,\text{ripple}}} \uparrow$, $T_{\mathrm{m}} \uparrow\uparrow$, $V_{\mathrm{m}} = \text{Unbalanced}$, $v \approx v_{\mathrm{ref}}$ & Localized insulation breakdown creates shorted winding loop with severe transformer circulating currents \cite{Gultekin2022SignalFusion, Li2023Prognostics}. & $\text{THD}(I_{\mathrm{m}}) > \text{Thresh} \land \frac{\mathrm{d}T_{\mathrm{m}}}{\mathrm{d}t} > \text{Baseline}$ \\
\addlinespace
\textbf{Battery SOH Degradation} & $I_{\mathrm{m}} = \text{Nom}$, $v = \text{Nom}$, $V_{\mathrm{bat}} \downarrow\downarrow \text{ under acceleration}$ & Increase in cell internal resistance ($R_{\mathrm{int}}$) due to chemical aging causes excessive transient voltage sag \cite{Jaramillo2026Battery}. & $\frac{\Delta V_{\mathrm{bat}}}{\Delta I_{\mathrm{bat}}} = R_{\mathrm{int}} > 1.5 \cdot R_{\mathrm{int,\text{nom}}}$ \\
\addlinespace
\textbf{Inverter Phase Loss} & $I_{\phi A} = 0$, $I_{\phi B,C} \uparrow\uparrow$, $v \ll v_{\mathrm{ref}}$, $T_{\mathrm{inv}} \uparrow\uparrow$ & Open-switch failure in inverter MOSFET gate or phase wire disconnection destroys rotating magnetic field \cite{Pillai2024Taxonomy, Li2023Prognostics}. & $\min(I_\phi) = 0 \land \max(I_\phi) > I_{\mathrm{rated}}$ \\
\bottomrule
\end{tabular}
\end{table}
